\Article{29}
Bevor wir weitergehen, ist zu bemerken, dass die Kriterien für Primzahlen (positiv genommen, wenn sie die Form $4n+1$ haben, negativ, wenn sie die Form $4n+3$ haben) zur Entscheidung für alle übrigen Zahlen ausreichen, sofern man den Satz aus Artikel 7 und die Kriterien für $-1$ und $\pm2$ hinzunimmt. Sucht man zum Beispiel Kriterien für die Zahl $+3$, so gelten die in Artikel 25 angegebenen Kriterien, die sich auf $-3$ beziehen, auch für $+3$, wenn $\frac12b$ gerade ist. Dagegen sind die Sammlungen $A,B,C,D$ mit den Sammlungen $C,D,A,B$ zu vertauschen, wenn $\frac12b$ ungerade ist. Daraus ergeben sich folgende Regeln:


\SourceLeaf{111}%
\begin{center}$+3$ gehört\end{center}
{\small
\[
\begin{array}{c|l}
\substack{\text{zur}\\\text{Sammlung}}&\text{wenn}\\
A&\begin{aligned}b\equiv0\ (\mathrm{mod.}\ 12);\\\text{oder zugleich }a\equiv0\ (\mathrm{mod.}\ 3),\ b\equiv2\ (\mathrm{mod.}\ 4)\end{aligned}\\
B&b\equiv8a\text{ oder }10a\ (\mathrm{mod.}\ 12)\\
C&\begin{aligned}b\equiv6a\ (\mathrm{mod.}\ 12);\\\text{oder zugleich }a\equiv0\ (\mathrm{mod.}\ 3),\ b\equiv0\ (\mathrm{mod.}\ 4)\end{aligned}\\
D&b\equiv2a\text{ oder }4a\ (\mathrm{mod.}\ 12)
\end{array}
\]
}

Ebenso erhält man die Kriterien für $\pm6$ durch Verbindung der Kriterien für $\mp2$ und $-3$, nämlich:
\begin{center}$+6$ gehört\end{center}
{\small
\[
\begin{array}{c|l}
\substack{\text{zur}\\\text{Sammlung}}&\text{wenn}\\
A&\begin{aligned}b\equiv0,\quad2a,22a\ (\mathrm{mod.}\ 24);\\\text{oder zugleich }a\equiv0\ (\mathrm{mod.}\ 3),\ b\equiv4a\ (\mathrm{mod.}\ 8)\end{aligned}\\
B&\begin{aligned}b\equiv4a,\quad6a,\quad8a\ (\mathrm{mod.}\ 24);\\\text{oder zugleich }a\equiv0\ (\mathrm{mod.}\ 3),\ b\equiv2a\ (\mathrm{mod.}\ 8)\end{aligned}\\
C&\begin{aligned}b\equiv10a,12a,14a\ (\mathrm{mod.}\ 24);\\\text{oder zugleich }a\equiv0\ (\mathrm{mod.}\ 3),\ b\equiv0\ (\mathrm{mod.}\ 8)\end{aligned}\\
D&\begin{aligned}b\equiv16a,18a,20a\ (\mathrm{mod.}\ 24);\\\text{oder zugleich }a\equiv0\ (\mathrm{mod.}\ 3),\ b\equiv6a\ (\mathrm{mod.}\ 8)\end{aligned}
\end{array}
\]
}

\begin{center}$-6$ dagegen gehört\end{center}
{\small
\[
\begin{array}{c|l}
\substack{\text{zur}\\\text{Sammlung}}&\text{wenn}\\
A&\begin{aligned}b\equiv0,\quad10a,14a\ (\mathrm{mod.}\ 24);\\\text{oder zugleich }a\equiv0\ (\mathrm{mod.}\ 3),\ b\equiv4a\ (\mathrm{mod.}\ 8)\end{aligned}\\
B&\begin{aligned}b\equiv4a,\quad8a,18a\ (\mathrm{mod.}\ 24);\\\text{oder zugleich }a\equiv0\ (\mathrm{mod.}\ 3),\ b\equiv6a\ (\mathrm{mod.}\ 8)\end{aligned}\\
C&\begin{aligned}b\equiv2a,12a,22a\ (\mathrm{mod.}\ 24);\\\text{oder zugleich }a\equiv0\ (\mathrm{mod.}\ 3),\ b\equiv0\ (\mathrm{mod.}\ 8)\end{aligned}\\
D&\begin{aligned}b\equiv6a,16a,20a\ (\mathrm{mod.}\ 24);\\\text{oder zugleich }a\equiv0\ (\mathrm{mod.}\ 3),\ b\equiv2a\ (\mathrm{mod.}\ 8)\end{aligned}
\end{array}
\]
}

Auf dieselbe Weise werden die Kriterien für die Zahl $+21$ aus den Kriterien für $-3$ und $-7$ zusammengestellt; die Kriterien für $-105$ aus denen für $-1$, $-3$, $+5$, $-7$ usw.

\Article{30}
Die induktive Untersuchung eröffnet somit eine überaus reiche Ernte besonderer Sätze, die mit dem Satz für die Zahl $2$ verwandt sind. Doch ein gemeinsames Band fehlt, strenge Beweise fehlen, da die Methode, mit der wir in der ersten Abhandlung die Behandlung der Zahl $2$ abgeschlossen haben, keine weitere Anwendung zulässt. Es gibt zwar verschiedene Methoden, mit denen sich Beweise für einzelne Fälle gewinnen ließen, insbesondere für diejenigen, die die Verteilung der quadratischen Reste auf die Sammlungen $A,C$ betreffen. Wir halten uns jedoch nicht bei ihnen auf, da eine allgemeine%
\SourceLeaf{112}%
, \emph{alle} Fälle umfassende Theorie unser Ziel sein muss. Nachdem wir im Jahr $1805$ begonnen hatten, dieser Frage unsere Überlegungen zu widmen, gelangten wir bald zu der Überzeugung, dass die eigentliche Quelle der allgemeinen Theorie in einer Erweiterung des Bereichs der Arithmetik zu suchen sei, wie wir bereits in Artikel 1 angedeutet haben.

Während sich die höhere Arithmetik bei den bisher behandelten Fragen allein mit reellen ganzen Zahlen befasst, erstrahlen die Sätze über biquadratische Reste erst dann in ihrer ganzen Einfachheit und eigentlichen Schönheit, wenn der Bereich der Arithmetik auf \emph{imaginäre} Größen ausgedehnt wird. Ohne Einschränkung bilden dann Zahlen der Form $a+bi$ ihren Gegenstand, wobei $i$ wie üblich die imaginäre Größe $\surd-1$ bezeichnet und $a,b$ uneingeschränkt alle reellen ganzen Zahlen zwischen $-\infty$ und $+\infty$ durchlaufen. Solche Zahlen werden wir \emph{komplexe ganze Zahlen} nennen; dabei werden die reellen Zahlen als eine darin enthaltene Art aufgefasst und den komplexen Zahlen nicht gegenübergestellt. Die vorliegende Abhandlung wird sowohl die elementare Lehre von den komplexen Zahlen als auch die ersten Anfänge der Theorie der biquadratischen Reste darstellen, deren allseitige Vollendung wir in einer nachfolgenden Fortsetzung unternehmen werden\begingroup\renewcommand{\thefootnote}{*)}\footnote{Wenigstens beiläufig ist hier noch zu bemerken, dass sich der so bestimmte Bereich besonders für die Theorie der biquadratischen Reste eignet. Die Theorie der kubischen Reste muss entsprechend auf die Betrachtung von Zahlen der Form $a+bh$ gegründet werden, wobei $h$ eine imaginäre Wurzel der Gleichung $h^3-1=0$ ist, etwa $h=-\frac12+\surd\frac34\pdot i$; ebenso wird die Theorie der Reste höherer Potenzen die Einführung anderer imaginärer Größen erfordern.}\endgroup.

\clearpage
\Article{31}
Zunächst führen wir einige Benennungen ein, die eine kürzere und klarere Darstellung ermöglichen.

Der Bereich der komplexen Zahlen $a+bi$ enthält
\begin{list}{}{\setlength{\leftmargin}{1.6em}\setlength{\labelwidth}{1.4em}\setlength{\labelsep}{.2em}}
\item[I.] die reellen Zahlen, bei denen $b=0$ gilt, und darunter je nach der Beschaffenheit von $a$
\begin{list}{}{\setlength{\leftmargin}{1.7em}\setlength{\labelwidth}{1.5em}\setlength{\labelsep}{.2em}}
\item[1)] die Null,
\item[2)] die positiven Zahlen,
\item[3)] die negativen Zahlen;
\end{list}
\item[II.] die imaginären Zahlen, bei denen $b$ von null verschieden ist. Hier unterscheiden wir wiederum
\begin{list}{}{\setlength{\leftmargin}{1.7em}\setlength{\labelwidth}{1.5em}\setlength{\labelsep}{.2em}}
\item[1)] imaginäre Zahlen ohne reellen Teil, also solche, bei denen $a=0$ gilt;
\item[2)] imaginäre Zahlen mit reellem Teil, wobei $b$ nicht null ist und auch $a=0$ nicht gilt.
\end{list}
\end{list}
Die ersteren kann man, wenn man will, rein imaginäre Zahlen, die letzteren gemischt imaginäre Zahlen nennen.

\SourceLeaf{113}%
In dieser Theorie verwenden wir vier Einheiten, $+1$, $-1$, $+i$, $-i$, die einfach positiv, negativ, positiv imaginär und negativ imaginär heißen sollen.

Die drei Produkte einer beliebigen komplexen Zahl mit $-1$, $+i$, $-i$ nennen wir ihre \emph{Gefährten} oder \emph{die ihr assoziierten Zahlen}. Mit Ausnahme der Null (die sich selbst assoziiert ist) sind also stets vier \emph{verschiedene} Zahlen einander assoziiert.

Dagegen nennen wir die Zahl, die aus einer komplexen Zahl durch Vertauschung von $i$ mit $-i$ entsteht, die zu ihr \emph{konjugierte} Zahl. Unter den imaginären Zahlen sind demnach stets zwei \emph{verschiedene} Zahlen konjugiert, während die reellen Zahlen sich selbst konjugiert sind, sofern man diese Benennung auf sie ausdehnen will.

Das Produkt einer komplexen Zahl mit der zu ihr konjugierten Zahl nennen wir die \emph{Norm} jeder der beiden. Bei einer reellen Zahl ist daher ihr Quadrat als ihre Norm anzusehen.

\Needspace{12\baselineskip}
Im Allgemeinen haben wir acht zusammenhängende Zahlen, nämlich
\[
\begin{array}{r|r}
a+bi&a-bi\\
-b+ai&-b-ai\\
-a-bi&-a+bi\\
b-ai&b+ai
\end{array}
\]
wobei wir zwei Gruppen von vier assoziierten Zahlen und vier Paare konjugierter Zahlen sehen; die gemeinsame Norm aller ist $aa+bb$. Die acht Zahlen reduzieren sich jedoch auf vier verschiedene, wenn entweder $a=\pm b$ oder für eine der beiden Zahlen $a,b=0$ gilt.

Aus den angeführten Definitionen ergibt sich unmittelbar Folgendes:

Zum Produkt zweier komplexer Zahlen ist das Produkt der zu ihnen konjugierten Zahlen konjugiert.

Dasselbe gilt für ein Produkt aus mehr Faktoren sowie für Quotienten.

Die Norm des Produkts zweier komplexer Zahlen ist gleich dem Produkt ihrer Normen.

Auch dieser Satz erstreckt sich auf Produkte aus beliebig vielen Faktoren und auf Quotienten.

Die Norm jeder komplexen Zahl (mit Ausnahme der Null, die wir von nun an meist stillschweigend ausschließen) ist eine \emph{positive} Zahl.

Ferner steht nichts im Wege, unsere Definitionen auf gebrochene oder sogar irrationale Werte von $a,b$ auszudehnen; doch soll $a+bi$ nur dann eine komplexe ganze Zahl heißen, wenn \emph{jede} der Zahlen $a,b$ ganz ist, und nur dann rational, wenn \emph{jede} der Zahlen $a,b$ rational ist.

